\documentclass[10pt,a4paper]{amsart}
\usepackage[margin=19mm]{geometry}
\usepackage{amsmath,amssymb,mathtools}
\usepackage[scr=rsfs,cal=euler]{mathalfa}
\usepackage{xfrac}
\usepackage[unicode,hidelinks,bookmarks=false]{hyperref}
\newcommand{\ie}{i.e.\ }
\newcommand{\cf}{cf.\ }
\newcommand{\resp}{respectively\ }
\newcommand{\Subsection}{\subsection}
\providecommand{\nobreakdash}{\nobreak\hbox{-}}
\setlength{\emergencystretch}{2em}
\multlinegap0pt
\usepackage{amssymb}
\usepackage{mathtools}
\usepackage[shortlabels]{enumitem}
\usepackage{xcolor}
\usepackage[all]{xy}
\usepackage{caption}
\usepackage{float}
\usepackage{csquotes}
\newtheorem{lemma}{Lemma}
\newtheorem{proposition}{Proposition}
\newcommand{\R}{\mathbb R}
\newcommand{\cT}[1]{\mathscr{T}_{#1}^{\scriptscriptstyle{>0}}}
\newcommand{\diam}{\mathrm{diam}}
\newcommand{\fF}{\mathfrak{F}}
\newcommand{\fS}{\mathfrak{S}}
\newcommand{\tL}{\mathtt{L}}
\newcommand{\tN}{\mathtt{N}}
\newcommand{\tP}{\mathtt{P}}
\title{Modular interpretation of the Weil--Petersson metric asymptotics for abelian varieties}
\author{Yanbo Fang, Andres Gomez}
\date{Source: arXiv:2602.23140v2 (CC BY 4.0)}
\begin{document}
\maketitle

\noindent\emph{Source and licence.} Modular interpretation of the Weil--Petersson metric asymptotics for abelian varieties. Version arXiv:2602.23140v2 (CC BY 4.0), distributed by arXiv under the Creative Commons Attribution 4.0 International licence, http://creativecommons.org/licenses/by/4.0/, as recorded in the arXiv licence metadata for this exact version and shown on the abstract page https://arxiv.org/abs/2602.23140v2. Full licence notice: \texttt{licenses/arxiv-2602.23140-CC-BY-4.0.txt}.\par\smallskip

\noindent\textit{Editorial scope.} Counted unit: the article's proposition pro:fiber.diam.est (source0.tex lines 1135--1149). The article's complete original proof of it is delivered with it (source0.tex lines 1150--1230, one \texttt{proof} environment of 81 lines). No mathematics is written, repaired or re-derived: the statement and the proof are the article's own lines, in the article's own notation, and no retained line is substituted or re-flowed.  Retained uncounted as the source-local context the proof needs: lines 1081--1093.  Nothing was omitted from any retained span, so the package carries no omission record.  No retained line contains a citation, so the wrapper carries no bibliography and no bibliography entry.  No retained line contains a figure environment, an image-inclusion command or drawing code, so no image is delivered and the package declares no asset dependency.  Editorial formatting only, in addition: an AMS \texttt{amsart} preamble that restates the article's own declarations and macros used by the retained lines, namely the environments \texttt{lemma}, \texttt{proposition} on separate counters, together with the standard packages those lines need.  The retained environments print in this document as Lemma 1, Proposition 1 in order of appearance, whereas the article prints the same items with the numbering of its own full document.  The complete native source of the article as distributed in the arXiv e-print for this exact version is bundled unchanged as \texttt{sources/195374/source0.tex}.
\begin{lemma}
\label{lem:vert.norm.est}
    Under the hypothesis of the previous lemma, there exists a constant $C>0$, depending only on the Siegel set data and $\tau'$, such that for every vertical tangent vector $V\in T_{\tau}\fF_g(u)$,
    \begin{equation}
        \label{eq:norm.est}
        \|V\|_{\tau}^2
        \le
        \frac{C}{\lambda_{\min}(t)}\Bigl(\|V_X'''\|_{\mathcal{F}}^2+\|V_Y'''\|_{\mathcal{F}}^2\Bigr)
        +
        \frac{C}{\lambda_{\min}(t)^2}\,\|V_X''\|_{\mathcal{F}}^2;
    \end{equation}
    here $\pi(\tau)=(\tau',t)$, $\lambda_{min}(t)$ denotes the smallest eigenvalue of $t$, and  by $\|\cdot\|_{\mathcal{F}}$ we denote the Frobenius norm.
\end{lemma}

\begin{proposition}\label{pro:fiber.diam.est}
    Let $\tau_n$ be a sequence in $\fF_g(u)\subset \fS_g$ such that $\tau_n\to \tau_\infty'\in \fF_{g'}(u)\subset \fS_{g'}$. This naturally defines a sequence in $\tP_{g'}\backslash\fS_g$, that we keep denoting by $\tau_n$. Let
    \[
        \pi_{g'}(\tau_n)=\big(\tau'_n\,,t_n\big)\in A_{g'}\times\cT{g''}.
    \]
    Then there exists a constant $C>0$ (depending only on the Siegel set and on $\tau'_\infty$) such that
    \begin{equation}\label{eq:diam-lambda}
        \diam\big(\pi_{g'}^{-1}\big(\tau'_n\,,t_n\big)\big) \le\ \frac{C}{\sqrt{\lambda_{\min}(t_n)}}.
    \end{equation}
    Furthermore, there exists $C'>0$ such that
    \begin{equation}\label{eq:diam-dgplus1}
        \diam\big(\pi_{g'}^{-1}\big(\tau'_n\,,t_n\big)\big)\leq \frac{C'}{\sqrt{d_{g'+1}(\tau_n)}}.
    \end{equation}
    In particular, as $\tau_n\to \tau'_\infty$ the diameter of the compact fiber vanishes.
\end{proposition}

\begin{proof}
Without loss of generality we work with the Siegel set description of the quotients, i.e.,
\[
    \pi_{g}:\fF_g(u)\longrightarrow\fF_{g'}(u)\times\fF^{+}_{g''}(u).
\]
Fix $(\tau',t)\in \fF_{g'}(u)\times\fF^{+}_{g''}(u)$ and let 
\[
    (X'''+iY''',\,X''),
    \qquad\text{so that}\qquad
    \tau''=X'' + i\bigl(t+{}^tY'''(Y')^{-1}Y'''\bigr).
\]
be the usual coordinates on the fiber over $(\tau',t)$. 

It is clear from the definiton of Siegel set that $\pi^{-1}_{g}(\tau',t)\cap \fF_g(u)$  gives a lattice fundamental domain of uniformly bounded size for the compact nilmanifold, possibly with finite orbifold structure, appearing as the fiber of the fibration\footnote{Recall that given a point $(\tau',t)$ in the base of the fibration $\pi_g: \tP_{g'}\backslash \fS_g\to A_{g'}\times\cT{g''}$ the fiber over it is homeomorpich to the (orbifold) compact nimanifold given by $\tL_{g',(\tau',t)}\backslash(\tN_{g'}\backslash N_{g'})$, where $\tL_{g',(\tau',t)}$ is the isotropy group of $(\tau',t)$}. Thus, there exists a constant $R>0$ (depending only on the chosen Siegel set) such that for any two points $\tau_p,\tau_q\in \pi_{g'}^{-1}(\tau',t)\cap\fF_g(u)$ the fiber-coordinate differences are bounded by $R$ in Frobenius norm, namely
\begin{equation}\label{eq:bounded-increments}
\|\Delta X'''\|_{\mathcal {F}}\le R,\qquad
\|\Delta Y'''\|_{\mathcal {F}}\le R,\qquad
\|\Delta X''\|_{\mathcal {F}}\le R.
\end{equation}

Given such $\tau_p,\tau_q$, connect them by a piecewise smooth path $\gamma(s)$ contained in the fiber, chosen so that $(\tau',t)$ remains constant and $X'''$ varies linearly on the first segment and, respectively, $Y'''$ and $X''$ on the second and third. Then $\dot\gamma(s)$ is vertical, and along each linear segment
\[
\|V_X'''\|_{\mathcal{F}}=\|\Delta X'''\|_{\mathcal{F}},\qquad
\|V_Y'''\|_{\mathcal{F}}=\|\Delta Y'''\|_{\mathcal{F}},\qquad
\|V_X''\|_{\mathcal{F}}=\|\Delta X''\|_{\mathcal{F}},
\]
and consequently
\[
\|V_X'''\|_{\mathcal{F}},\;\|V_Y'''\|_{\mathcal{F}},\;\|V_X''\|_{\mathcal{F}}\;\le R.
\]
Applying Lemma~\ref{lem:vert.norm.est} to $V=\dot\gamma(s)$ yields
\[
\|\dot\gamma(s)\|^2
\le
\frac{C_0}{\lambda_{\min}(t)}\Bigl(\|V_X'''\|_{\mathcal{F}}^2+\|V_Y'''\|_{\mathcal{F}}^2\Bigr)
+
\frac{C_0}{\lambda_{\min}(t)^2}\,\|V_X''\|_{\mathcal{F}}^2
\le
\frac{C_1}{\lambda_{\min}(t)},
\]
for a constant $C_1$ depending only on the Siegel set and $\tau'$. In particular, $\|\dot\gamma(s)\|\le C_2/\sqrt{\lambda_{\min}(t)}$. Integrating along $\gamma$ gives
\[
d(\tau_p,\tau_q)\le \mathrm{length}(\gamma)
\le \frac{C_{\tau'}}{\sqrt{\lambda_{\min}(t)}},
\]
uniformely over all such $\tau_p,\tau_q\in \pi_{g'}^{-1}(\tau',t)\cap \fF_g(u)$, with $C_{\tau'}$ depending only on the Siegel set and $\tau'$. Taking the supremum over $\tau_p,\tau_q$ yields 
\[
\diam\big(\pi_{g'}^{-1}\big(\tau'\,,t\big)\big) \le\ \frac{C_{\tau'}}{\sqrt{\lambda_{\min}(t)}}.
\]
Then, applying this to the sequence $(\tau'_n,t_n)$ and noticing that $\tau'_n\to \tau_\infty$ we get \eqref{eq:diam-lambda}.

Finally, \eqref{eq:diam-dgplus1} follows from the fact that, given a point $(\tau',t)$ and any $\tau\in \pi^{-1}_{g'}(\tau',t)$, the Siegel set inequalities $d_{i}~\leq~ u \,d_{i+1}$ imply the existence of $c>0$ such that
\begin{equation}\label{eq:lambdamin-lower}
    \lambda_{\min}(t)\ \ge\ c\, d_{g'+1}(\tau).
\end{equation}
To see this, write the Jacobi's decomposition $Y=LDL^t$ in block form as
\begin{equation}
       \begin{pmatrix}
            L'&0\\
            L'''&L''
        \end{pmatrix}
        \begin{pmatrix}
            D'&0\\
            0&D''
        \end{pmatrix}
        \begin{pmatrix}
            ^tL'&^tL'''\\
            0&^tL''
        \end{pmatrix}.
\end{equation}
This gives $t=L''\,D''\,{}^tL''$ with $D''=\mathrm{diag}(d_{g'+1},\cdots,d_{g})$. Let $x \in \R^{g''}$. Then
\[
{}^tx\, t\, x = {}^t\bigl({}^tL'' x\bigr)\,D''\,\bigl({}^tL'' x\bigr) \ge d_{min}\,\|{}^tL'' x\|^2.
\]
Using $\|{}^tL'' x\|\ge \frac{\|x\|}{\|(L'')^{-1}\|_{\mathrm{op}}}$ and taking the infimum over $\|x\|=1$ yields
\[
\lambda_{min}(t) \ge \frac{d_{min}}{\|(L'')^{-1}\|_{\mathrm{op}}^2},
\]
%The expression above follows from the inequality $\lambda_{min}(t) \|x\| \leq {}^tx\,t\,x \leq \lambda_{max}(t)\|x\|$
where $d_{min}=\mathrm{min}\{d_i\,|\, g'+1\leq i \leq g \}$. Thus \eqref{eq:lambdamin-lower} follows from $\|(L'')\|_{\mathrm{op}}$ being uniformly bounded on the Siegel set (and consequently $\|(L'')^{-1}\|_{\mathrm{op}}$), together with $d_{min} \geq d_{g'+1}/\mathrm{max}\{1,u^{g''-1}\}$.
\end{proof}

\end{document}
